\section{Deferral and Boundary}

The downstream evidence boundary specification reserves RAUDIT and RCERT but
does not implement them. The civic-evidence access-pattern specification gives
the read direction, and the RCOUNT audit-algorithm roadmap shows why the current
audit work still belongs inside RCOUNT \citep{civic-evidence-boundaries-2026}.
The V-series contracts are still proving the common transcript shape: method
ids, assertion ids, sample or batch references, risk limits, per-step statistics,
final decisions, and explicit boundary outcomes.

RAUDIT should therefore stay embedded until a second production consumer needs
the same audit transcript independently of an RCOUNT package. That consumer
might be a public case bundle, a cross-jurisdiction audit library, or a
certification record that wants to cite a transcript package by hash. Until that
pressure exists, extracting RAUDIT would create another package boundary without
making the evidence more verifiable.

RCERT has a different deferral reason. Certification is source-artifact driven.
Before schema work begins, the project needs real public certification
artifacts: board minutes, canvass reports, motions, official certificates,
signatory records, evidence checklists, legal caveats, and any referenced audit,
log, or custody exhibits. Without those artifacts, RCERT would risk inventing a
generic workflow rather than preserving what officials actually publish.

One decision remains open: whether the V-series audit-algorithm transcripts need
standalone RAUDIT packages at all. Some methods may remain most useful as
embedded RCOUNT replay contracts. Others may become reusable exhibits that a
certification or public-case layer should reference directly. The boundary rule
does not answer that packaging question. It only fixes ownership: RAUDIT can
replay method evidence. RCERT can record official actions. Neither owns the base
count ledger, and neither should silently change RCOUNT verifier results.
